% Partie : logique intuitionniste
% Chapitre : propositions comme types
% Section : des termes aux preuves

\documentclass[../../../include/open-logic-section]{subfiles}

\begin{document}

\olfileid[fr]{int}{pty}{tp}

\olsection{Reconstruire des \usetoken{p}{derivation} à partir des termes de preuve}

Considérons le sens inverse : reconstruire à partir d'un terme
de preuve~$N$ une preuve dans \Log{N2i}.
On fixe un contexte contenant une affectation $x:!A$
pour chaque variable libre~$x$ de~$N$. On vérifie ensuite
que $N$ est bien typé dans ce contexte.
Le typage donne l'arbre de la dérivation annotée.
Commençons cependant par un exemple sans variable libre :
\[
\lambd[\typeof{x}{!A \land
  !B}][\pair{\proj{2}{x}}{\proj{1}{x}}]
\]
Ce terme est de la forme $\lambd[\typeof{x}{!A \land !B}][N_1]$.
La seule règle de \Log{tN2i} qui produit un terme de cette
forme est~\Intro{\lif}. La preuve encodée par~$N$
se termine donc par cette règle :
\[
  \Axiom$x : !A\land !B \fCenter \pair{\proj{2}{x}}{\proj{1}{x}} : 
    !C$
  \RightLabel{\Intro\lif}
  \UnaryInf$\fCenter \lambd[\typeof{x}{!A \land
  !B}][\pair{\proj{2}{x}}{\proj{1}{x}}] : (!A \land !B) \lif !C$
  \DisplayProof
  \]
Nous ne connaissons pas encore $!C$, mais l'antécédent
de l'implication finale est $!A \land !B$, et le contexte
de la prémisse doit contenir $x:!A \land !B$.

Le terme associé à cette prémisse est déjà déterminé :
$\pair{\proj{2}{x}}{\proj{1}{x}}$.
La preuve qu'il encode se termine nécessairement par
\Intro{\land}. La formule $!C$ est donc une conjonction,
$!C \ident (!C_1 \land !C_2)$.
La reconstruction se poursuit :
\[
  \Axiom$x : !A\land !B \fCenter \proj{2}{x} : !C_1$
  \Axiom$x : !A\land !B \fCenter \proj{1}{x} : !C_2$
  \RightLabel{\Intro\land}
  \BinaryInf$x : !A\land !B \fCenter \pair{\proj{2}{x}}{\proj{1}{x}} : 
    !C_1 \land !C_2$
  \RightLabel{\Intro\lif}
  \UnaryInf$\fCenter \lambd[\typeof{x}{!A \land
  !B}][\pair{\proj{2}{x}}{\proj{1}{x}}] : (!A \land !B) \lif (!C_1 \land !C_2)$
  \DisplayProof
  \]
Le terme de preuve $\proj{2}{x}$ indique que le séquent
supérieur gauche est obtenu par \Elim{\land}[2] :
sa prémisse a donc la forme $!D \land !C_1$.
À droite, l'inférence qui produit $!C_2$ est
\Elim{\land}[1], et sa prémisse a la forme
$!C_2 \land !E$ :
\[
  \Axiom$x : !A\land !B \fCenter x: !D\land !C_1$
  \RightLabel{\Elim\land[2]}
  \UnaryInf$x : !A\land !B \fCenter \proj{2}{x} : !C_1$
  \Axiom$x : !A\land !B \fCenter x: !C_2\land !E$
  \RightLabel{\Elim\land[1]}
  \UnaryInf$x : !A\land !B \fCenter \proj{1}{x} : !C_2$
  \RightLabel{\Intro\land}
  \BinaryInf$x : !A\land !B \fCenter \pair{\proj{2}{x}}{\proj{1}{x}} : 
    !C_1 \land !C_2$
  \RightLabel{\Intro\lif}
  \UnaryInf$\fCenter \lambd[\typeof{x}{!A \land
  !B}][\pair{\proj{2}{x}}{\proj{1}{x}}] : (!A \land !B) \lif (!C_1 \land !C_2)$
  \DisplayProof
  \]
Les termes des séquents supérieurs sont maintenant des
variables : ces séquents sont des axiomes.
On en déduit $!C_1$, $!C_2$, $!D$ et $!E$.
À gauche, $!D \ident !A$ et $!C_1 \ident !B$,
sans quoi le séquent ne serait pas un axiome.
À droite, $!C_2 \ident !A$ et $!E \ident !B$
pour la même raison. La preuve est entièrement reconstruite :
\[
  \Axiom$x : !A\land !B \fCenter x: !A\land !B$
  \RightLabel{\Elim\land[2]}
  \UnaryInf$x : !A\land !B \fCenter \proj{2}{x} : !B$
  \Axiom$x : !A\land !B \fCenter x: !A\land !B$
  \RightLabel{\Elim\land[1]}
  \UnaryInf$x : !A\land !B \fCenter \proj{1}{x} : !A$
  \RightLabel{\Intro\land}
  \BinaryInf$x : !A\land !B \fCenter \pair{\proj{2}{x}}{\proj{1}{x}} : 
    !B \land !A$
  \RightLabel{\Intro\lif}
  \UnaryInf$\fCenter \lambd[\typeof{x}{!A \land
  !B}][\pair{\proj{2}{x}}{\proj{1}{x}}] : (!A \land !B) \lif (!B \land !A)$
  \DisplayProof
  \]
En effaçant les termes de preuve des conséquents, on retrouve la dérivation de \Log{N2i}.
\end{document}
